% Generated draft: requires mathematical and source review.
\documentclass[11pt]{article}
\usepackage[margin=25mm]{geometry}
\usepackage{fontspec}
\setmainfont{Latin Modern Roman}
\usepackage{amsmath,amssymb,mathtools,unicode-math}
\setmathfont{Latin Modern Math}
\usepackage{enumitem,needspace,xurl,hyperref}
\setlength{\parindent}{0pt}
\setlength{\parskip}{5pt}
\setlength{\emergencystretch}{3em}
\newcommand{\N}{\mathbb{N}}
\newcommand{\Z}{\mathbb{Z}}
\newcommand{\Q}{\mathbb{Q}}
\newcommand{\R}{\mathbb{R}}
\newcommand{\C}{\mathbb{C}}
\newcommand{\F}{\mathbb{F}}
\newcommand{\A}{\mathbb{A}}
\newcommand{\PP}{\mathbb{P}}
\newcommand{\E}{\mathbb{E}}
\newcommand{\eps}{\varepsilon}
\newcommand{\dd}{\,\mathrm{d}}
\begin{document}
\section*{UnsolvedMath ID: 30003408}
\textbf{Source problem code:} catalogue entry
\subsection{Definitions and mathematical statement}
Let \(G\) be a topological group. A \(G\)-flow is a nonempty compact Hausdorff space \(X\) equipped with a jointly continuous action \(G\times X\to X\). It is minimal if every orbit is dense. A universal minimal flow \(M(G)\) is a minimal \(G\)-flow such that every minimal \(G\)-flow \(X\) admits a continuous surjective \(G\)-equivariant map \(M(G)\to X\). Metrizability means that the topology of \(M(G)\) is induced by a metric.

A subgroup \(H\leq G\), with its subspace topology, is extremely amenable if every \(H\)-flow has a point fixed by every element of \(H\). It is co-precompact in \(G\) if, for every identity neighbourhood \(U\subseteq G\), there is a finite set \(F\subseteq G\) such that \(G=UFH\).

\textbf{General criterion problem.} Find a verifiable condition \(\mathcal C\) on \(G\), together with a proof of
\[
\mathcal C(G)\ \Longrightarrow\ M(G)\text{ is metrizable},
\]
that is easier to check than the existence of a closed, co-precompact, extremely amenable subgroup \(H\leq G\). The source gives no formal meaning to ``easier to check.'' It also gives no explicit restriction of the ambient class of topological groups, such as a Polish-group hypothesis, and does not specify separation axioms for \(G\). These omissions require review: the comparison must not be interpreted as asserting that the subgroup condition alone implies metrizability for every topological group.

\textbf{Particular question in the supplied expanded statement.} Let \(L\) be a finite relational language, meaning that it has finitely many relation symbols, each of finite arity. Let \(\mathcal A\) be a countable \(L\)-structure that is homogeneous: every isomorphism between finite induced substructures of \(\mathcal A\) extends to an automorphism of \(\mathcal A\). Equip \(G=\operatorname{Aut}(\mathcal A)\) with the topology of pointwise convergence on the underlying discrete set. Thus a neighbourhood basis at the identity consists of the subgroups
\[
G_{(B)}=\{g\in G: g(b)=b\text{ for every }b\in B\},
\qquad B\subseteq A\text{ finite},
\]
where \(A\) is the underlying set of \(\mathcal A\). Is \(M(\operatorname{Aut}(\mathcal A))\) metrizable for every such \(\mathcal A\)? The pointwise-convergence topology and the stated meaning of homogeneity are conventional interpretations of terminology left implicit in the supplied statement. The record does not specify whether ``countable'' means at most countable or countably infinite.
\subsection{Short English statement}
Find a practical sufficient condition ensuring that a topological group's universal minimal flow is metrizable, with a precise account of its applicable hypotheses. The expanded source statement also asks whether this always holds for automorphism groups of countable homogeneous structures in finite relational languages, using the conventional pointwise-convergence topology. The ambient hypotheses and the meaning of an easier criterion require clarification.
\subsection{Sources}
\begin{itemize}
\item Local UnsolvedMath registry, supplied complete source record: permanent ID \(30003408\), exact source problem code \texttt{catalogue entry}, title ``Practical Criteria for Metrizable Universal Minimal Flows,'' in the Oberwolfach Reports Open Problems collection. The record supplies proposed year \(2017\) and identifies no individual proposer.
\item Supplied report citation: \emph{Set Theory (2018)}, Oberwolfach Reports, pp.~527--589, DOI \texttt{10.4171/owr/2017/11}. \url{https://doi.org/10.4171/owr/2017/11}. The supplied citation's year is retained despite its mismatch with the DOI year and proposed year.
\item Supplied literature source, ``OWR Set Theory report'': \url{https://ems.press/content/serial-article-files/46674}.
\item Supplied literature source, ``Metrizable universal minimal flows'': \url{https://academic.oup.com/imrn/article/2016/5/1285/2450868}.
\item Literature assessment dated \textbf{2026-08-21}: the supplied curation labels the problem open and states that the report leaves an easier general criterion open, while later work characterizes metrizable universal minimal flows without supplying such a criterion in full generality. This is supplied dated curation, not an independently verified literature assessment; its time-sensitive status and cited claims require independent verification.
\item The supplied record labels the statement \texttt{corrected\_verified} and reports a direct check of the official report. That verification claim is preserved as source metadata; no independent source consultation or mathematical verification is claimed here.
\end{itemize}
\end{document}
